\section{Raising/lowering indexes}

At this point, we're pretty familiar with vector spaces and dual vector spaces.\\
But we can ask if, there's a way to create a correspondence between the vectors of $\textit{V}$ and the covectors of $V^*$?
In other words, is there any meaningfull way that we can take a vector in $\textit{V}$ and find its partner in $V^*$?\\

We already saw one way, where we paired up the basis vectors and basis covectors by choosing the basis covectors acting on the basis vectors
so that they map as Kroneker delta.
Right off the bat, since this works in any dimention, we can think that, given a vector $\vec{v}$ in a specific basis $\vec{e_i}$, to find its 
partner in the dual space, we just expand it as a linear combination of the basis vectors (imagine dim=4):

\begin{equation}
    \begingroup\color{orange}\vec{v}\endgroup = \begingroup\color{cyan}v^1\endgroup \begingroup\color{blue}\vec{e_1}\endgroup + 
    \begingroup\color{cyan}v^2\endgroup \begingroup\color{blue}\vec{e_2}\endgroup + \begingroup\color{cyan}v^3\endgroup \begingroup\color{blue}\vec{e_3}\endgroup +
    \begingroup\color{cyan}v^4\endgroup \begingroup\color{blue}\vec{e_4}\endgroup
\end{equation}

We then turn the basis vectors into basis covectors and we can say that, the partner covector of vector $\begingroup\color{orange}\vec{v}\endgroup$ is:

\begin{equation}
    \begingroup\color{orange}\alpha\endgroup = \begingroup\color{cyan}v^1\endgroup \begingroup\color{blue}\epsilon^1\endgroup + 
    \begingroup\color{cyan}v^2\endgroup \begingroup\color{blue}\epsilon^2\endgroup + \begingroup\color{cyan}v^3\endgroup \begingroup\color{blue}\epsilon^3\endgroup +
    \begingroup\color{cyan}v^4\endgroup \begingroup\color{blue}\epsilon^4\endgroup
\end{equation}

This seems very nice and straight forward at first, but there's a problem with this approach, and this is very clear
from a simple example.\\
Let's consider a change of basis, where we multiply the basis vectors by a scaling factor of 2:

\begin{equation}
    \begingroup\color{red}\tilde{\vec{e_i}}\endgroup = 2 \begingroup\color{blue}\vec{e_i}\endgroup
\end{equation}

So the forward and backward transforms, are very simply:

\begin{align*}
    F = \begin{bmatrix}
        2 & 0\\
        0 & 2
        \end{bmatrix}  &&  B = \begin{bmatrix}
        1/2 & 0\\
        0 & 1/2
        \end{bmatrix}
\end{align*}

The problem here, is that basis vectors are covariant, so they transform from the old 
to the new, using the forward transform, but basis covectors are contravariant, so they transform, from
old to new, using the backward transform.\\
So this method of aligning partners do not work very well, because if the correspondence may look nice in one basis, 
it starts to look weird and wrong when we perform a basis/coordinate change.

So we need to try better, and maybe this time without using basis at all.\\
What we propose trying is, taken a vector $\begingroup\color{orange}\vec{v}\endgroup$, we introduce the covector \textit{v-dot-something}
$\left(\begingroup\color{orange}\vec{v}\endgroup \cdot  {{} - {}} \right)$ that can take another vector as argument 
and outputs a scalar.\\
We may wonder if this covector is indeed part of the dual space.
Well first of all, this is for sure a function from our vector space $\textit{V}$ to a scalar, because
if we supply a vector, we know that the dot product produces a number.

\begin{equation}
    \left(\begingroup\color{orange}\vec{v}\endgroup \cdot  {{} - {}} \right) : V \to \mathbb{R}
\end{equation}

This has to be linear in order to be part of the dual space, right?
But this is indeed, due to properties of the dot product, so eventually:

\begin{equation}
    \left(\begingroup\color{orange}\vec{v}\endgroup \cdot  {{} - {}} \right) \in V^*
\end{equation}

Differently than before, now we see that we're not dependent anymore to the basis vectors we choose.
If before, growing the basis vectors meant shrinking the basis covectors, now scaling up or down the vector means also
scaling by the same amount the correspondent covector \textit{v-dot-something}\\

Not, if this covector lives in the dual space, that means we should be able to write it 
as linear combination of basis covectors, so something like:

\begin{equation}
    \begingroup\color{orange}\vec{v}\endgroup \cdot  {{} - {}} = x_i \begingroup\color{blue}\epsilon^i\endgroup
\end{equation}

What are the $x_i$ coefficients?\\
We recall that the result of the dot product of two vectors is given by passing the vectors to the metric
tensor. 

\begin{align*}
    \begingroup\color{orange}\vec{v}\endgroup \cdot \begingroup\color{violet}\vec{w}\endgroup = \begingroup\color{cyan}g_{ij}\endgroup \begingroup\color{orange}v^i\endgroup \begingroup\color{violet}w^j\endgroup
\end{align*}

In this case now, we can do:

\begin{align*}
    \begingroup\color{orange}\vec{v}\endgroup \cdot  {{} - {}} &= 
    g\left(\begingroup\color{orange}\vec{v}\endgroup, -\right)\\ 
    &= \begingroup\color{cyan}g_{ik}\endgroup \begingroup\color{blue}\epsilon^i\epsilon^k\endgroup \left(\begingroup\color{orange}v^j\endgroup \begingroup\color{blue}\vec{e_j}\endgroup\right)\\ 
    &= \begingroup\color{cyan}g_{ik}\endgroup \begingroup\color{orange}v^j\endgroup \begingroup\color{blue}\epsilon^i\epsilon^k\endgroup \left( \begingroup\color{blue}\vec{e_j}\endgroup\right)
\end{align*}

Since the metric tensor is symmetric, we can arbitrarily choose to pass the basis vector to either one of the two covectors there, producing 
a Kroneker delta:

\begin{align*}
    \begingroup\color{orange}\vec{v}\endgroup \cdot  {{} - {}} &= 
    \begingroup\color{cyan}g_{ik}\endgroup \begingroup\color{orange}v^j\endgroup \begingroup\color{blue}\epsilon^i\epsilon^k\endgroup \left( \begingroup\color{blue}\vec{e_j}\endgroup\right)\\
    &= \begingroup\color{cyan}g_{ik}\endgroup \begingroup\color{orange}v^j\endgroup \begingroup\color{blue}\epsilon^i\endgroup \delta^k_j\\
    &= \begingroup\color{cyan}g_{ij}\endgroup \begingroup\color{orange}v^j\endgroup \begingroup\color{blue}\epsilon^i\endgroup
\end{align*}

And we see that $\begingroup\color{cyan}g_{ij}\endgroup \begingroup\color{orange}v^j\endgroup$ are indeed the components of the 
covector \textit{v-dot-something} in the epsilon (dual vector) basis.\\

Now, as an alternative notation, we can also say that:

\begin{align*}
    \begingroup\color{orange}\vec{v}\endgroup \cdot  {{} - {}} = g\left(\begingroup\color{orange}\vec{v}\endgroup, -\right) &= \begingroup\color{cyan}g_{ij}\endgroup \begingroup\color{orange}v^j\endgroup \begingroup\color{blue}\epsilon^i\endgroup\\
    \begingroup\color{orange}\vec{v}\endgroup \cdot  {{} - {}} = g\left(\begingroup\color{orange}\vec{v}\endgroup, -\right) &= \begingroup\color{orange}v_j\endgroup \begingroup\color{blue}\epsilon^i\endgroup
\end{align*}

What this notation tells us, is that it's almost as if the metric tensor components are lowering the index of the components of $\vec{v}$ to 
give the covector components of \textit{v-dot-something}:

\begin{align*}
    \begingroup\color{orange}\vec{v}\endgroup = \begingroup\color{orange}v^j\endgroup \begingroup\color{blue}\vec{e_j}\endgroup &&
    \begingroup\color{orange}\vec{v}\endgroup \cdot  {{} - {}} = \begingroup\color{orange}v_j\endgroup \begingroup\color{blue}\epsilon^i\endgroup
\end{align*}

In summary:

\begin{subequations}\label{eq:forward-backward}
\begin{empheq}[box=\widefbox]{align}
    \begingroup\color{orange}\vec{v}\endgroup &= \begingroup\color{orange}v^j\endgroup \begingroup\color{blue}\vec{e_j}\endgroup = \begingroup\color{brown}\tilde{v^j}\endgroup \begingroup\color{red}\tilde{\vec{e_j}}\endgroup\\
    \begingroup\color{orange}\vec{v}\endgroup \cdot  {{} - {}} &= \begingroup\color{orange}v_i\endgroup \begingroup\color{blue}\epsilon^i\endgroup = \begingroup\color{brown}\tilde{v_i}\endgroup \begingroup\color{red}\tilde{\epsilon^i}\endgroup\\
    \begingroup\color{orange}v_i\endgroup &= \begingroup\color{cyan}g_{ij}\endgroup \begingroup\color{orange}v^j\endgroup\\ 
    \begingroup\color{brown}\tilde{v_i}\endgroup &= \begingroup\color{magenta}\tilde{g_{ij}}\endgroup \begingroup\color{brown}\tilde{v^j}\endgroup
\end{empheq}
\end{subequations}

What's really important to note down, is that:

\begin{equation}
    \begingroup\color{orange}v_i\endgroup \neq \begingroup\color{orange}v^i\endgroup
\end{equation}

The vector components are not the same as the partner covector components, and to switch between them, 
you have to use the metric tensor components and do a summation.
The only case where these components are the same, is when the metric tensor components are given by the 
Kroneker delta, which basically means that we're working with an orthonormal coordinate system.\\

So, we found a way to partner up vectors and covectors, by pairing the vector with the covector \textit{v-dot-something}, and 
the way we get from the vector to covector is by using the metric tensor.
We normally think of the metric tensor being a function from a pair of vectors to scalars, but
it can be also thought as a function from a vector space to its dual:

\begin{align*}
    g : V \times V \to \mathbb{R}\\
    g : V \to V^*
\end{align*}

But what about the reverse direction? How do we get from a covectors to its paired vector partner?\\
Well we know all about the metric tensor, we know its components have to lower indexes and we know that
it's a member of the tensor product:

\begin{equation}
    g \in V^* \otimes V^*
\end{equation}

We're now going to introduce what's known as the "inverse metric tensor":

\begin{equation}
    \mathfrak{g} \in V \otimes V
\end{equation}

By definition, this is defined so that, when combined with the ordinary metric tensor, it produces the Kroneker delta:

\begin{equation}
    \begingroup\color{blue}\mathfrak{g}^{ki}\endgroup \begingroup\color{cyan}g_{ij}\endgroup = \delta^k_i
\end{equation}

We can try that out by simply using the definition:

\begin{align*}
    \begingroup\color{orange}v_i\endgroup &= \begingroup\color{cyan}g_{ij}\endgroup \begingroup\color{orange}v^j\endgroup\\
    \begingroup\color{blue}\mathfrak{g}^{ki}\endgroup \begingroup\color{orange}v_i\endgroup &= \begingroup\color{blue}\mathfrak{g}^{ki}\endgroup \begingroup\color{cyan}g_{ij}\endgroup \begingroup\color{orange}v^j\endgroup\\
    \begingroup\color{blue}\mathfrak{g}^{ki}\endgroup \begingroup\color{orange}v_i\endgroup &= \delta^k_j \begingroup\color{orange}v^j\endgroup\\
    \begingroup\color{blue}\mathfrak{g}^{ki}\endgroup \begingroup\color{orange}v_i\endgroup &= \begingroup\color{orange}v^k\endgroup
\end{align*}

If ordinary metric tensor lower indexes, the inverse metric tensor goes in the other direction, so it raises indexes.\\
To give a summary image:

\begin{center}
\input{figures/08-from-v-to-dual-metric}
\end{center}

The ordinary metric tensos is also referred as the covariant metric tensor (lower indexec), and the inverse
metric tensor is also referred as contravariant metric tensor.\\

Now, this raising/lowering indexes does not only apply to vectors and covectors components, but to other tensors too.
Just an example, let's say we have a tensor Q:

\begin{align*}
    Q &\in V \otimes V^* \otimes V^*\\
    Q &= Q^i_{jk} \begingroup\color{blue}\vec{e_i}\epsilon^j\epsilon^k\endgroup
\end{align*}

If we multiply by the inverse metric tensor $\mathfrak{g}^{jk}$, we raise the j index:

\begin{equation}
    Q^i_{jk} \begingroup\color{blue}\mathfrak{g}^{jx}\endgroup = Q^{ix}_k
\end{equation}

This new tensor obtained $Q'$ is a member of $V \otimes V \otimes V^*$:

\begin{align*}
    Q' &\in V \otimes V \otimes V^*\\
    Q' &= Q^{ix}_k \begingroup\color{blue}\vec{e_i}\vec{e_x}\epsilon^k\endgroup
\end{align*}